\chapter{Introduction}\label{ch:intro}

This note accompanies the paper on relativistic detuning, electron heating and the reach of continuum kinetic methods in resonant dark-photon conversion in a plasma. The paper is a short report of results. The note is the long version of everything the paper takes for granted: the field theory of the dark photon, the kinetic plasma theory that is used to judge the numerical results, the numerical methods and the way they are put together in Dark-SPECTRAX. It is meant for a reader who knows classical electrodynamics and some kinetic theory and wants to follow the argument without going through twenty references first. The chapters are independent enough to be read in isolation; Chapters~\ref{ch:darkphoton} to~\ref{ch:numerics} are self-contained derivations, and the later chapters use them to discuss the results.

\section{Why this problem}

Dark matter has not been detected through any non-gravitational interaction, and the search has widened from weakly interacting particles of GeV mass to very light bosons, which behave as classical fields when their occupation numbers are large. One of the simplest candidates is a massive vector boson, a dark photon, that mixes kinetically with the ordinary photon. Its Lagrangian has one extra gauge field, one mass and one dimensionless mixing parameter $\chi$, and because the mixing is a renormalizable operator it can be generated by heavy particles charged under both gauge groups at loop level~\citep{Holdom1986}. The review of Fabbrichesi, Gabrielli and Lanfranchi collects the theory and the laboratory and astrophysical searches~\citep{Fabbrichesi2020}. If its mass lies below roughly an electron volt and a mechanism populates the field in the early Universe, the dark photon behaves as cold dark matter, with a coherent field oscillating at its Compton frequency~\citep{NelsonScholtz2011,Arias2012}.

A coherent dark-photon field does not interact with a vacuum, but it does with a plasma. The kinetic mixing makes the dark field exert a force on charges that is $\eta\simeq\chi$ times the force of an ordinary field of the same strength (Sec.~\ref{sec:mixing}). In a plasma the electrons oscillate collectively at the plasma frequency $\wpe$, and when $\wpe$ equals the dark-photon mass the dark field drives this oscillation resonantly, in the way a swing is driven by a push at its natural frequency. In the early Universe the free electrons and protons that survive recombination form a plasma whose density, and hence $\wpe$, falls with expansion. For a dark photon of any mass in a wide range there is a moment at which the plasma frequency sweeps through the mass. Energy then moves from the dark field into plasma oscillations, which damp and heat the gas. The effect is the plasma analogue of the conversion of dark photons into photons in a medium, and its Landau--Zener description is the subject of Sec.~\ref{sec:resonance}~\citep{Mirizzi2009,DubovskyHC2015,McDermottWitte2020}.

\section{The bounds and the claim of Hook, Huang and Shalaby}

Heating of the intergalactic gas and distortions of the cosmic microwave background by conversion at the resonance put limits on the mixing as a function of the dark-photon mass. The dark-photon limits are summarized in a recent handbook~\citep{CaputoHandbook2021}, and the cosmological ones depend on how much of the dark-matter energy is deposited in the plasma when the resonance is crossed~\citep{McDermottWitte2020,Witte2020}. The calculation assumes that the transfer is described by linear response: the dark field drives a longitudinal plasma oscillation, the oscillation is damped by collisions or by Landau damping, and the energy ends up as heat. In this picture the deposited energy is a fixed fraction of the dark-matter energy density that depends on the Landau--Zener adiabaticity parameter of Eq.~\eqref{eq:LZ}, and nothing limits the amplitude of the plasma oscillation until the resonance is passed.

Hook, Huang and Shalaby (HHS) argued that this picture fails because of plasma nonlinearity~\citep{HHS2026}. Their reasoning is that at resonance the amplitude of the $k=0$ Langmuir oscillation grows secularly, and once the energy density of the wave approaches the thermal energy of the plasma, the ponderomotive force of the wave pushes the electrons and ions into density inhomogeneities. The plasma frequency becomes a function of position, the resonance is detuned in each region, and conversion stops. They support this with one-dimensional relativistic particle-in-cell simulations, with a prescribed spatially uniform dark-photon drive, mass ratio $m_i/m_e=1836$, and drive strengths $v_q/v_{te}=0.03$ and $10^{-3}$, where $v_q=E_0/\wpe$ is the velocity scale of the drive (Sec.~\ref{sec:resonance}), and with Landau--Zener sweeps. The energy deposited is of the order of the electron thermal energy at the time of conversion, which they find to be orders of magnitude below the amount assumed in the cosmological bounds, so that the constraints weaken by factors between $3\times10^3$ and $10^7$ over ten decades in mass. The result was published as Phys.\ Rev.\ Lett.\ 137, 071004. It is a strong claim: if it holds, it removes a class of published bounds, and it rests on a numerical experiment.

The main conclusion of the claim is physical and the supporting evidence is numerical, and the numerical method has a feature that matters. A particle-in-cell simulation represents the plasma by a finite number of particles and it is noisy; the noise is not a defect to be removed, because it is also a property of the real plasma, in which there are thermal fluctuations with the same spectrum, only at a lower relative amplitude owing to the number of particles. If the instabilities that stop the transfer are seeded by noise, then the saturated state depends on the noise level, which in a PIC code is set by the number of particles per cell, and the question of which number is the physically correct one is not trivial.

\section{What a noise-free kinetic method can and cannot settle}

A continuum kinetic method represents the distribution function itself. It has no particle noise, so each run is deterministic, the work done by the drive and the energy of the fields and of the plasma are known with an energy ledger that closes to the tolerance of the integrator, and results at different resolutions can be compared without averaging over seeds. The method of this note represents the distribution in velocity by Hermite functions and in space by Fourier modes (Chapter~\ref{ch:numerics}). It can compute the linear response to the accuracy of the discretization and compare it with the exact kinetic roots, and it can run a prescribed resonant drive and measure when the response departs from the linear law. With it we can decide whether a given departure is due to a process present in the equations (as opposed to noise), by turning the process off: removing the ions, fixing the temperature, switching between nonrelativistic and relativistic kinematics, or changing the basis. Because the right-hand side is a smooth function of the model parameters, the same method can in principle be differentiated.

It also has limits, which determine what can be said about HHS. A spectral method is not positivity-preserving: the truncated expansion can give a negative distribution function once structures finer than the resolution develop. For the asymmetrically weighted basis it has in addition a numerical instability that bounds the length of a reliable run; the paper defines certified windows by declared checks and reports only what is inside them. The method cannot represent cell-scale shot noise as a physical fluctuation, because with PIC-like particle numbers its distribution function loses positivity almost immediately. Three-dimensional velocity space and magnetized plasmas are available in the code, but the production runs are one-dimensional; collisions are included only in the nonrelativistic model and with a velocity-independent frequency. The method is not a substitute for cosmological modelling: it knows nothing of expansion, cooling by Compton scattering or the temperature-dependent collision frequency.

With these capabilities the questions that the method can settle are the following. Where, in a model with the physical ingredients of HHS, does the energy transfer from a resonant uniform drive leave the linear resonant law? What sets the timing of that departure? What is the heating that HHS plot, and can a deterministic calculation produce it? The answers of the paper, which the later chapters derive and qualify, are short. The timing is that of relativistic detuning, the decrease of the plasma frequency of the quivering electrons as their mass grows: both relativistic continuum models leave the linear law by 10\,\% at $\wpe t\simeq392.5$ for $\vq=0.03$, and the homogeneous relativistic Vlasov equation at $k=0$, a few lines of code, gives 392.3 (Sec.~\ref{sec:relativistic}). The heating is a different matter. An independent relativistic PIC converges to about 25 initial thermal energies at $\wpe t=1000$, half the 50 read off the figure of HHS, a Newtonian pusher gives 63, and the deterministic continuum codes give $0.056$, a factor of 445 below the PIC. The origin of the factor of two between the PIC and HHS is not resolved, and whether finite-amplitude fluctuations absent from the continuum seeds produce the heating is an open question. We do not claim that HHS is wrong.

The nonrelativistic results that preceded the relativistic ones are retained with their scope marked. They concern a model whose coherent quiver speed reaches $0.47$--$0.87c$ at the ends of its certified windows (1000, 892 and 791.5 at $\vq=0.03$, 0.05 and 0.07), so they are statements about that model and not about a relativistic plasma (Chapter~\ref{ch:nonrel}).

\section{Roadmap}

The remainder of the first part of the note has three chapters. Chapter~\ref{ch:darkphoton} derives the dark-photon model from the kinetic-mixing Lagrangian and its coupling to a plasma, including the resonance and the Landau--Zener passage. Chapter~\ref{ch:plasma} collects the kinetic plasma theory against which the numerical results are judged, from Landau damping to the parametric instabilities of a uniform pump that are central to the HHS argument, and ends with the relativistic oscillator. Chapter~\ref{ch:numerics} describes the numerical methods, particle-in-cell, grid, discontinuous Galerkin and Hermite, and the way they appear in Dark-SPECTRAX. Chapters~\ref{ch:model} to~\ref{ch:limits} then specify the model that is solved, the ladder of verification and the limits of the Hermite solver, and Chapters~\ref{ch:hhs} to~\ref{ch:rel} take up the HHS claim and the nonrelativistic and relativistic studies. Chapter~\ref{ch:discussion} collects what is settled and what is open, and Appendix~\ref{app:repro} lists repositories, pins and records.

Units are $c=\wpe=n=1$ unless stated, with $\vq=v_q/v_{te}$ for the drive strength. Numbers quoted from the repository carry their record or section in a parenthesis or footnote, with paths relative to the repository root. Where the note computes something that is not in a record, it says so and gives the script (\texttt{note/make\_note\_figures.py}). Literature is cited where it is used, and only for statements that the cited work contains.
